\section{Introduction}
\label{sec:introduction}

\subsection{Three Levels of Redistricting Perturbation}

Modern algorithmic redistricting operates across a hierarchy of perturbation
scales.
Each level explores a different radius of the plan space, incurs a different
computational cost, and serves a different analytical purpose.

\paragraph{Level 1: Plan-level (global) perturbation.}
Each invocation of the METIS partitioner with a different random seed produces
an independent redistricting plan.
These plans are unrelated in the graph-partition sense; there is no continuous
path between them in the space of valid plans.
\cs{}~\citep{b11paper} runs $T = 600$ independent seeds and selects the
minimum-edge-cut plan; \textsc{PercentileSweep} selects the plan at a given
percentile of the edge-cut distribution.
Cost: $O(T \cdot n \log n / k)$ per state.
Exploration radius: global (the full feasible plan space is sampled, subject
to the METIS prior).

\paragraph{Level 2: Node-level (intermediate) perturbation.}
\be{}~\citep{h1paper} runs short \recom{} chains at each node of the
bisection tree.
These chains explore the neighbourhood of a fixed hierarchical decomposition,
resampling sub-district assignments while holding the higher levels fixed.
Cost: $O(\text{nodes} \times \text{chain steps} \times n/k)$.
Exploration radius: moderate (the hierarchical structure constrains reachable
plans, but within each node the local plan space is sampled).

\paragraph{Level 3: Tract-level (local) perturbation.}
The \flip{} chain moves a single boundary tract from its current district to
an adjacent district at each step.
This is the coarsest grain at which a valid redistricting plan can be perturbed
while remaining a connected-district plan.
Cost: $O(1)$ per step (one adjacency query, one contiguity check).
Exploration radius: local (the chain is confined to the neighbourhood of
the initial plan; see Section~\ref{sec:formal}).

This paper studies Level 3 in depth.
We argue that \flip{} is unsuitable for ensemble sampling of the full plan
space---the mixing time is far too large---but is the right tool for
\emph{local sensitivity analysis}: given a submitted plan, how stable are
its partisan outcomes under single-tract perturbations?

\subsection{Why Local Sensitivity Analysis Matters}

A redistricting plan submitted to a court or legislature is not a point
estimate; it is a political act.
Opponents will seek the nearest plan that produces a different partisan
outcome---and if such a plan exists within a single tract move, the submitted
plan is arguably unstable.

Local sensitivity analysis under \flip{} provides a principled quantification
of this instability:
\begin{enumerate}
  \item Run a \flip{} chain from the submitted plan.
  \item Record all visited plans and their partisan seat outcomes.
  \item Report the \emph{partisan stability} --- the fraction of visited plans
        with the same seat outcome as the submitted plan.
\end{enumerate}
A plan with partisan stability $\geq 95\%$ is robust to single-tract
perturbations; a plan with stability $< 80\%$ lies near a partisan boundary
in the plan space, which is potential evidence of intentional placement.

\subsection{Main Contributions}

\begin{enumerate}
  \item \textbf{Algorithm.} We give the complete specification of \flip{} as
        implemented in \texttt{redist-cli} (\texttt{SeedCompositor::Flip}),
        including the deterministic seed schedule and the acceptance criterion
        (Section~\ref{sec:formal}).

  \item \textbf{Formal properties.} We prove that (a) each step runs in $O(1)$
        time excluding the contiguity BFS (which is $O(n/k)$ in the worst case
        but $O(\deg)$ on average), (b) the acceptance rate is high ($\approx
        40$--$70\%$), (c) the chain mixes slowly (exponential mixing time in
        the worst case), and (d) not all valid plans are reachable from a
        given start via \flip{} moves (Section~\ref{sec:formal}).

  \item \textbf{Sensitivity metrics.} We define EC stability and partisan
        stability, and demonstrate their behaviour on North Carolina
        ($k = 14$) and Wisconsin ($k = 8$) (Section~\ref{sec:sensitivity}).

  \item \textbf{Comparison.} We provide a unified comparison of \flip{},
        \sburst{}, \be{}, and \cs{} across cost, acceptance rate, exploration
        radius, and use case (Section~\ref{sec:comparison}).

  \item \textbf{Audit chain.} We describe the manifest record for \flip{}
        runs and show that recording \texttt{visited\_count} and
        \texttt{selected\_plan\_rank} provides a reproducibility guarantee
        stronger than seed-only recording (Section~\ref{sec:audit}).
\end{enumerate}

\subsection{Paper Structure}

Section~\ref{sec:formal} defines the \flip{} algorithm formally and derives
its complexity and mixing properties.
Section~\ref{sec:sensitivity} introduces the sensitivity metrics and reports
empirical results for NC and WI.
Section~\ref{sec:comparison} compares \flip{} with ReCom-based methods.
Section~\ref{sec:audit} describes the audit chain and reproducibility protocol.
Section~\ref{sec:conclusion} concludes with usage guidelines.
